%%==========================================================================
\section{The Maintenance-Extended Model and Correction of Constraint~(13)}
\label{sec:maintenance}
%%==========================================================================

\subsection{Maintenance problem definition}

Following \citet{khaled2018compact}, we consider type-A checks triggered after
at most $T^A_{\max}$ cumulative flight hours.  Checks take place at night
after the last flight on a given day in an airport with maintenance capacity.

Additional notation is collected in \Cref{tab:maint_notation}.

\begin{table}[t]
\centering\small
\caption{Additional notation for the maintenance model.}
\label{tab:maint_notation}
\begin{tabular}{lp{7cm}}
\toprule
Symbol & Meaning \\
\midrule
$\D = \{1,\ldots,n_d\}$      & Planning days \\
$d_i$                         & Departure day of flight $i$ \\
$\MA \subseteq \A$            & Maintenance-capable airports \\
$\text{mcap}_{m,d}$           & Maintenance capacity of airport $m$ on day $d$ \\
$mc_{ijd}$                    & Cost of night maintenance at arrival airport of
                                flight $i$, aircraft $j$, day $d$ \\
$F(d,i)$                      & $\{i' : \tD{i'} > \tI{i},\; d_{i'} = d_i\}$ \\
$T^A_{\max}$                  & Max cumulative flight minutes before type-A check \\
$d_{\max}$                    & Max calendar-day interval between checks \\
$M_{dd'}$                     & Big-$M$ for day pair $(d,d')$ \\
$F_{d,d'}$                    & Flights departing strictly between days $d$ and $d'$ \\
\bottomrule
\end{tabular}
\end{table}

\subsection{Additional decision variables}

\begin{align*}
  z_{ijd} &= \begin{cases}1 & \text{night maintenance at arrival airport of } i,\;
               \text{aircraft } j,\;\text{day } d, \\ 0 & \text{else;}\end{cases}\\
  y_{jd}  &= \begin{cases}1 & \text{maintenance of aircraft } j \text{ on day } d,
               \\ 0 & \text{else.}\end{cases}
\end{align*}

\subsection{Extended model}

The objective becomes:
\begin{equation}\label{eq:obj_ext}
  \min\; \sum_{i \in \F}\sum_{j \in \Pc} c_{ij} x_{ij}
  + \sum_{m \in \MA}\sum_{i \in \F^{\downarrow}(m)}\sum_{j \in \Pc}\sum_{d \in \D}
    mc_{ijd}\, z_{ijd}.
\end{equation}

Constraints~\eqref{eq:c1}--\eqref{eq:c3} are retained.  Additional
constraints:

\paragraph{C8 — Maintenance blocks same-day subsequent flights.}
\begin{equation}\label{eq:c8}
  z_{ijd} + x_{i'j} \le 1,
  \qquad \forall i \in \F,\; i' \in F(d,i),\; j \in \Pc,\; d \in \D.
\end{equation}

\paragraph{C9 — Maintenance requires assignment.}
\begin{equation}\label{eq:c9}
  x_{ij} \ge z_{ijd},
  \qquad \forall i \in \F,\; j \in \Pc,\; d \in \D.
\end{equation}

\paragraph{C10 — Capacity.}
\begin{equation}\label{eq:c10}
  \sum_{i \in \F^{\downarrow}(m)} \sum_{j \in \Pc} z_{ijd}
  \le \text{mcap}_{m,d},
  \qquad \forall d \in \D,\; m \in \MA.
\end{equation}

\paragraph{C11 — Aggregation ($z \to y$).}
\begin{equation}\label{eq:c11}
  \sum_{i \in \F} z_{ijd} = y_{jd},
  \qquad \forall d \in \D,\; j \in \Pc.
\end{equation}

\paragraph{C12 — Maximum-spacing constraint.}
\begin{equation}\label{eq:c12}
  \sum_{r \in [d, d']} y_{jr} \ge 1,
  \qquad \forall d, d' \in \D \text{ s.t.\ } d' - d = d_{\max} - 1,\;
  \forall j \in \Pc.
\end{equation}

\paragraph{C14 — At most one maintenance event per aircraft per day.}
\begin{equation}\label{eq:c14}
  \sum_{i \in \F} z_{ijd} \le 1,
  \qquad \forall j \in \Pc,\; d \in \D.
\end{equation}

\subsection{Impact of corrected C2--C4 feasibility on C8--C14}
\label{subsec:c2c4_impact_maintenance}

The correction introduced in \Cref{sec:basic} (non-overlap
constraint~\eqref{eq:c4_overlap}) changes the admissible assignment set for
$x$, and therefore the interpretation of all maintenance constraints layered on
top of $x$.

\paragraph{C8.}
No algebraic change is required in~\eqref{eq:c8}; however, with
\eqref{eq:c4_overlap} active, the blocking relation now operates on physically
realizable trajectories only. This removes spurious cases where maintenance
placement was evaluated on infeasible simultaneous departures.

\paragraph{C9--C12 and C14.}
Constraints~\eqref{eq:c9}--\eqref{eq:c12}, \eqref{eq:c14} remain structurally unchanged. Their
feasible region is tightened indirectly because $z$ and $y$ are linked to a
strictly smaller, physically valid $x$-space.

\paragraph{C13 and onward.}
Hour-accumulation constraints are now evaluated along executable flight paths;
this avoids artificial inflation/deflation of cumulative hours caused by
simultaneous assignments that cannot occur in operations.

In summary, the C2--C4 correction propagates to C8--C14 primarily as a
\emph{domain correction}: the maintenance constraints keep their form, but are
enforced over a physically consistent assignment polytope.

\subsection{Rigorous feasibility theorem for C1--C14}
\label{subsec:feasibility_theorem_c1_c14}

We now state the corrected sufficiency result for the type-A maintenance model.
Consider the constraint family
\eqref{eq:c1}--\eqref{eq:c3}, \eqref{eq:c4_overlap}, \eqref{eq:c8}--\eqref{eq:c12},
\eqref{eq:c13_new_a}--\eqref{eq:c13_new_b}, \eqref{eq:c14}, and binary domains.

\begin{theorem}[Physical feasibility under corrected C1--C14]
  \label{thm:physical_feasibility_c1_c14}
  Any binary solution satisfying the above constraints induces, for each
  aircraft $j$, a physically executable trajectory over the planning horizon:
  flights assigned to $j$ are temporally non-overlapping, satisfy airport
  continuity and turn time, and admit a day-by-day maintenance schedule
  satisfying capacity and hour-spacing rules.
\end{theorem}

\begin{proof}
  Fix an aircraft $j$ and let
  $S_j=\{i\in\F: x_{ij}=1\}$ ordered by nondecreasing departure time.

  \textbf{(i) Flight assignment and continuity.}
  By~\eqref{eq:c1}, every flight is assigned to exactly one aircraft.
  By~\eqref{eq:c2}--\eqref{eq:c3}, each departure assigned to $j$ has sufficient
  prior balance at its departure airport, with the initial unit at $a^0_j$.
  Therefore airport continuity holds along $S_j$.

  \textbf{(ii) Temporal executability.}
  For any overlapping pair $(i_1,i_2)\in\mathcal{O}$,
  \eqref{eq:c4_overlap} enforces $x_{i_1j}+x_{i_2j}\le 1$.
  Hence no two flights assigned to $j$ overlap once turn time is included.
  Combined with (i), this yields a realizable single-aircraft flight timeline.

  \textbf{(iii) Maintenance event consistency.}
  By~\eqref{eq:c9}, maintenance trigger variables can only be activated on
  assigned flights. By~\eqref{eq:c11}, $y_{jd}=1$ iff exactly one daily trigger
  is selected in aggregate; by~\eqref{eq:c14}, at most one maintenance event is
  attached to aircraft $j$ on day $d$.

  \textbf{(iv) Maintenance capacity and blocking.}
  By~\eqref{eq:c10}, airport-day capacities are respected. By~\eqref{eq:c8}, if
  maintenance is triggered after flight $i$ on day $d$, then any subsequent
  same-day flight is forbidden for the same aircraft, so the end-of-day
  maintenance action is operationally executable.

  \textbf{(v) Hour/day maintenance admissibility.}
  By~\eqref{eq:c12}, checks cannot be spaced farther than $d_{\max}$ days.
  By corrected constraints~\eqref{eq:c13_new_a}--\eqref{eq:c13_new_b},
  cumulative flight time between relevant maintenance delimiters is bounded by
  $T^A_{\max}$ (cf. \Cref{thm:c13_valid}).

  Steps (i)--(v) show that every aircraft has a consistent, non-overlapping
  flight sequence with feasible maintenance placement and resource usage.
  Therefore the global solution is physically feasible.
\end{proof}

\subsection{Constraint~(13): original formulation and its flaw}
\label{subsec:c13_flaw}

The original paper's cumulative-flight-time constraint is:
\begin{equation}\label{eq:c13_paper}
  \tag{13$_{\mathrm{paper}}$}
  \sum_{i \in F_{d,d'}} x_{ij}\,\tilde{t}_i
  \;\le\;
  T^A_{\max}
  + M_{dd'}\Bigl(2 - y_{jd} - y_{jd'}\Bigr)
  + M_{dd'}\!\sum_{r \in [d+1,d'-1]} y_{jr},
\end{equation}
for all $j \in \Pc$, all pairs $(d, d')$ with $2 \le d'-d \le d_{\max}-1$,
$d < n_d$, where $\tilde{t}_i$ is the flight duration of leg $i$.

\begin{lemma}[Flaw in~\eqref{eq:c13_paper}]
  \label{lem:c13_flaw}
  Consider two consecutive maintenance events at days $d$ and $d'$ with no
  check between them ($y_{jr} = 0$ for $r \in (d, d')$).  Then the RHS
  of~\eqref{eq:c13_paper} equals $T^A_{\max} + M_{dd'}(2 - 1 - 1) +
  M_{dd'} \cdot 0 = T^A_{\max}$, which correctly imposes the flight-hour
  limit between the two checks.

  However, if $y_{jd} = 0$ (no maintenance at the start day) and $y_{jd'} =
  1$ (maintenance only at the end day), the RHS becomes $T^A_{\max} +
  M_{dd'}(2 - 0 - 1) = T^A_{\max} + M_{dd'}$, which is vacuously large and
  imposes \emph{no limit} on cumulative flight hours before the first check of
  the period.  Hence the constraint fails to enforce the check-interval bound
  for the opening sub-period.
\end{lemma}

\begin{proof}
  Direct substitution of $y_{jd} = 0$, $y_{jd'} = 1$, $\sum_{r} y_{jr} = 0$
  into~\eqref{eq:c13_paper} gives the stated RHS.  Since $M_{dd'}$ is chosen
  as the maximum possible cumulative flight time over the horizon (a large
  positive constant), the resulting constraint is trivially satisfied regardless
  of the actual flight hours flown.
\end{proof}

\subsection{Corrected Constraint~(13)}
\label{subsec:c13_corrected}

We propose replacing~\eqref{eq:c13_paper} with two separate constraints, each
anchored to one endpoint:

\begin{equation}\label{eq:c13_new_a}
  \tag{13a}
  \sum_{i \in F_{d,d'}} x_{ij}\,\tilde{t}_i
  \;\le\;
  T^A_{\max}
  + M_{dd'}\, y_{jd}
  + M_{dd'}\!\sum_{r \in [d+1,d'-1]} y_{jr},
\end{equation}
\begin{equation}\label{eq:c13_new_b}
  \tag{13b}
  \sum_{i \in F_{d,d'}} x_{ij}\,\tilde{t}_i
  \;\le\;
  T^A_{\max}
  + M_{dd'}\, y_{jd'}
  + M_{dd'}\!\sum_{r \in [d+1,d'-1]} y_{jr},
\end{equation}
for all $j \in \Pc$, $(d,d') \in \D^2$ with $2 \le d'-d \le d_{\max}-1$,
$d < n_d$.

\begin{theorem}[Validity of corrected Constraint~(13)]
  \label{thm:c13_valid}
  A feasible assignment satisfying \eqref{eq:c13_new_a}--\eqref{eq:c13_new_b}
  for all aircraft and all day pairs correctly enforces the cumulative
  flight-hour bound $T^A_{\max}$ between every pair of consecutive maintenance
  events, regardless of whether either endpoint day $d$ or $d'$ hosts a check.
\end{theorem}

\begin{proof}
  Let $d^* < d^{**}$ be two consecutive maintenance days for aircraft $j$
  (i.e., $y_{jd^*} = 1$, $y_{jd^{**}} = 1$, $y_{jr} = 0$ for $r \in
  (d^*, d^{**})$).  For the pair $(d, d') = (d^*, d^{**})$:
  \begin{itemize}
    \item In \eqref{eq:c13_new_a}: RHS $= T^A_{\max} + M_{dd'} \cdot 1 +
      M_{dd'} \cdot 0 = T^A_{\max} + M_{dd'}$, which does not constrain.
    \item In \eqref{eq:c13_new_b}: RHS $= T^A_{\max} + M_{dd'} \cdot 1 +
      M_{dd'} \cdot 0 = T^A_{\max} + M_{dd'}$, same.
  \end{itemize}
  Now consider a sub-interval $(d, d')$ where $y_{jd} = 0$ and $y_{jd'} = 0$
  (no check at either end):
  \begin{itemize}
    \item Both \eqref{eq:c13_new_a} and \eqref{eq:c13_new_b} give
      RHS $\le T^A_{\max} + M_{dd'}\sum_r y_{jr}$.  If no check exists in
      $(d, d')$ either ($\sum_r y_{jr} = 0$), both constraints impose LHS
      $\le T^A_{\max}$, correctly bounding the sub-period flight hours.
  \end{itemize}
  In all cases the constraints are active precisely when they should be,
  confirming validity.
\end{proof}

\begin{remark}[Tightness]
  The corrected formulation \eqref{eq:c13_new_a}--\eqref{eq:c13_new_b} is no
  weaker than \eqref{eq:c13_paper} (both are implied by the same underlying
  combinatorial condition), but is \emph{strictly tighter} when one endpoint
  indicator is zero: the original constraint admits an $M_{dd'}$ slack at one
  endpoint, which the corrected pair does not.  This yields smaller LP
  relaxation bounds and fewer branch-and-bound nodes (confirmed
  computationally in \Cref{sec:computational}).
\end{remark}

\subsection{Pre-horizon accumulated hours (Constraint~13b)}
\label{subsec:c13b}

The original paper does not account for maintenance hours accumulated before
the planning horizon begins.  We add a constraint for the opening sub-period
$[0, d']$:
\begin{equation}\label{eq:c13b}
  \tag{13c}
  \sum_{i \in F_{0,d'}} x_{ij}\,\tilde{t}_i
  \;\le\;
  \bigl(T^A_{\max} - \Phi^A_j\bigr)
  + M_{0d'}\, y_{jd'}
  + M_{0d'}\!\sum_{r \in [1, d'-1]} y_{jr},
\end{equation}
where $\Phi^A_j \ge 0$ is the pre-horizon accumulated flight time for aircraft
$j$ toward its next type-A check.  This constraint is only generated when
$\Phi^A_j > 0$.

\subsection{Integrality constraints}

\begin{equation}
  x_{ij} \in \{0,1\},\quad
  z_{ijd} \in \{0,1\},\quad
  y_{jd} \in \{0,1\},
  \qquad \forall i,j,d.
\end{equation}
